% Rapatrié depuis le doc de travail (onglet Article FR), révision 144.

\section{Discussion}
\label{sec:7}

\subsection{Apports théoriques}
\label{sec:7.1}

La première contribution est méthodologique. L'étude de la résilience pilotée bute sur l'absence de données reliant une perturbation, les décisions prises et la trajectoire de performance. La démarche proposée produit de telles données et leur donne une propriété que les données d'observation n'ont jamais : chaque crise y est observée sous plusieurs politiques, toutes choses égales par ailleurs. La politique devient ainsi une intervention au sens de \citet{pearl2009}, dont l'effet est identifiable par construction, et une préconisation peut être confrontée à ce qui se serait réellement produit.

La deuxième contribution porte sur la représentation de la résilience. Le triangle de résilience \citep{bruneau2003} et le profil de perturbation \citep{sheffi2005book} résument une crise par sa profondeur et sa durée. Les résultats montrent que ces deux dimensions ne relèvent pas des mêmes causes : la profondeur de la chute dépend surtout du choc, la gestion n'en modifie presque pas le minimum, alors que la durée de la crise dépend surtout de la gestion, le retour à la normale passant du quatorzième au quatrième jour selon la politique. Décomposer la trajectoire en paramètres de mécanisme, comme le fait la courbe hyperbolique composée, n'est donc pas un raffinement descriptif : c'est la condition pour attribuer une perte à sa cause. Le réseau appris retrouve d'ailleurs ce partage sans qu'il lui soit imposé.

La troisième contribution précise la typologie des stratégies d'atténuation et de contingence \citep{tomlin2006,chopra2004}. L'efficacité d'un levier n'est pas une propriété du levier, mais de sa rencontre avec la nature du choc et la topologie du réseau : le déstockage est sans effet sur une panne fournisseur mais atténue une congestion, l'efficacité de la réaction la plus rapide s'effondre lorsque les voies de secours sont étroites, et la couverture de stock protège contre les pertes de capacité de distribution sans protéger contre un choc de demande. Ces interactions plaident pour des règles de pilotage conditionnelles plutôt que génériques, que seul un modèle causal du réseau considéré permet d'établir.

\subsection{Implications managériales et économiques}
\label{sec:7.2}

Les indicateurs de résilience se traduisent directement en volume de ventes perdues. Sur le réseau ISOMORPH au niveau de tension étudié, une crise sans réaction coûte en médiane 1,34 jour de demande, soit environ 21 000 unités. Le profil réactif en évite en moyenne 1,06 jour, le profil prudent 0,34 et le profil tardif 0,22. Pour une marge unitaire m, la valeur d'une politique se lit donc comme le produit du service évité, de la demande journalière et de m, à comparer au coût des leviers qu'elle engage ; le modèle ne valorisant pas ce coût, cette comparaison revient au décideur, mais l'ordre de grandeur du gain est désormais connu.

Quatre implications managériales s'en dégagent. Premièrement, la vitesse de détection et de décision semble compter davantage que le choix des leviers : investir dans la surveillance de la performance et dans des délégations de décision pré-établies paraît prioritaire. Deuxièmement, l'information sur la durée prévisible d'un choc a une valeur économique propre : elle permet de gérer environ une crise sur sept avec une politique plus sobre, pour une perte de réussite inférieure à deux points, et la communication avec fournisseurs et transporteurs sur la durée des incidents mérite d'être organisée avant la crise. Troisièmement, les points de passage obligés réduisent l'efficacité de toute réaction et justifient des investissements de redondance ciblés, que les leviers de crise ne remplacent pas. Quatrièmement, les stocks de sécurité ne sont pas une protection universelle, et une réserve d'urgence gagne à être tenue séparée de la politique de commande, faute de quoi elle peut perturber le signal de réapprovisionnement (section~\ref{sec:3.5}).

\subsection{Limites et perspectives}
\label{sec:7.3}

Les résultats reposent sur des données simulées et en héritent les limites. Les valeurs posées faute de données de terrain, capacités, effets des leviers, caractéristiques des profils, ont été choisies pour être plausibles et sont toutes réglables, mais elles conditionnent les ordres de grandeur ; seules les tendances et les hiérarchies, stables entre les deux réseaux étudiés, doivent être généralisées. Le niveau de tension retenu et le relèvement des seuils de tolérance visent à produire des crises contrastées : ils ne décrivent pas une fréquence réelle de crises. Une validation sur des données d'entreprise, même partielles, reste nécessaire.

Plusieurs simplifications du gestionnaire simulé bornent la portée des préconisations. Les leviers n'ont pas de coût et restent engagés jusqu'à la fin de la période observée, ce qui rend la domination de la réactivité en partie structurelle ; introduire leur coût et leur levée transformerait la préconisation en arbitrage explicite entre coût et résilience. La nature de la perturbation est supposée diagnostiquée sans erreur, chaque scénario ne comporte qu'une perturbation, et la politique est représentée par le profil seul : en distinguer les composantes suppose des lots où elles varient indépendamment.

Le pronostic est limité par l'information disponible au moment de décider plutôt que par la taille de la base : la tension du réseau manque au contexte, et les crises touchant un point de passage obligé sont trop rares pour que leurs effets propres soient appris, ce qu'un échantillonnage stratifié corrigerait.

La représentation temporelle appelle enfin deux réserves. La couverture de stock est uniforme entre sites, et son interaction avec la politique de gestion n'a pas été étudiée. La recherche de la demande de référence par dichotomie suppose une monotonie qui n'est pas démontrée, même si la valeur retenue est toujours vérifiée par simulation.